\section{Chuẩn độ acid--base}

\subsection{Chuẩn độ dùng để làm gì?}

Chuẩn độ là phương pháp xác định lượng hoặc nồng độ của một chất bằng cách cho chất đó phản ứng với một dung dịch đã biết nồng độ.

Trong chuẩn độ acid--base, một dung dịch được thêm từ từ vào dung dịch kia cho đến khi phản ứng vừa đủ theo tỉ lệ phương trình hóa học.

\begin{ghinho}{Ba ý cần hiểu trước khi tính}
\begin{enumerate}
\item \textbf{Dung dịch chuẩn} có nồng độ đã biết.
\item \textbf{Điểm kết thúc} được nhận biết nhờ dấu hiệu thích hợp, thường là sự đổi màu của chất chỉ thị.
\item Thể tích dung dịch chuẩn đã dùng cho biết số mol chất chuẩn đã phản ứng; từ phương trình hóa học suy ra số mol chất cần tìm.
\end{enumerate}
\end{ghinho}

\subsection{Dụng cụ và cách đọc burette}

Burette cho phép thêm dung dịch từng lượng nhỏ và đo chính xác thể tích đã dùng. Thể tích dung dịch chảy ra được tính bằng:
\[
V_{\text{đã dùng}}=V_{\text{cuối}}-V_{\text{đầu}}.
\]

Khi gần điểm kết thúc, cần thêm dung dịch từng giọt và lắc đều bình tam giác. Nếu dùng phenolphthalein để chuẩn độ acid bằng base, màu hồng rất nhạt bền trong thời gian ngắn thường được dùng làm dấu hiệu kết thúc.

\begin{notechem}
Không được xem ``đổi màu'' là nguyên nhân của phản ứng. Chất chỉ thị chỉ giúp quan sát thời điểm phản ứng đã đạt gần tỉ lệ vừa đủ.
\end{notechem}

\subsection{Con đường giải bài chuẩn độ}

\begin{ghinho}{Hướng giải}
\[
\boxed{C_{\text{chuẩn}},V_{\text{chuẩn}}}
\to
\boxed{n_{\text{chuẩn}}}
\to
\boxed{\text{tỉ lệ PTHH}}
\to
\boxed{n_{\text{chất cần tìm}}}
\to
\boxed{C,\ m,\ \%\ldots}
\]
\end{ghinho}

\begin{vidu}
Lấy \SI{20.0}{\milli\liter} dung dịch \ce{HCl} chưa biết nồng độ, chuẩn độ vừa đủ bằng \SI{16.0}{\milli\liter} dung dịch \ce{NaOH} \SI{0.100}{\mole\per\liter}. Tính nồng độ mol của \ce{HCl}.
\end{vidu}
\begin{loigiai}
Số mol \ce{NaOH} đã dùng:
\[
n_{\ce{NaOH}}=0.0160\times0.100=\SI{0.00160}{\mole}.
\]
\[
\ce{HCl + NaOH -> NaCl + H2O}.
\]
Tỉ lệ mol là $1:1$ nên:
\[
n_{\ce{HCl}}=\SI{0.00160}{\mole}.
\]
Do đó:
\[
C_{\ce{HCl}}=\frac{0.00160}{0.0200}=\SI{0.0800}{\mole\per\liter}.
\]

\textbf{Vậy,} nồng độ mol của dung dịch \ce{HCl} là \SI{0.0800}{\mole\per\liter}.
\end{loigiai}

\subsection{Thực hành mô phỏng}

Hãy thực hiện ít nhất ba lần chuẩn độ và ghi kết quả vào bảng:

\begin{center}
\begin{tabular}{c|ccc}
\toprule
Lần & $V_{\text{đầu}}$ (mL) & $V_{\text{cuối}}$ (mL) & $V_{\text{đã dùng}}$ (mL)\\
\midrule
1 & & & \\
2 & & & \\
3 & & & \\
\bottomrule
\end{tabular}
\end{center}

Có thể dùng một trong các mô phỏng miễn phí sau:
\begin{itemize}
\item RSC Titration Screen Experiment: \href{https://edu.rsc.org/resources/titration-screen-experiment/2077.article}{edu.rsc.org/resources/titration-screen-experiment/2077.article}.
\item ChemCollective Virtual Lab: \href{https://chemcollective.org/vlab}{chemcollective.org/vlab}.
\end{itemize}

Sau khi thao tác, trả lời bốn câu hỏi: dung dịch nào đã biết nồng độ; dung dịch nào cần tìm; dấu hiệu nào cho biết đã gần điểm kết thúc; vì sao phải thêm từng giọt khi gần điểm kết thúc?

\begin{tuluyen}
Chuẩn độ \SI{10.0}{\milli\liter} dung dịch \ce{NaOH} bằng dung dịch \ce{HCl} \SI{0.0500}{\mole\per\liter}. Thể tích \ce{HCl} trung bình cần dùng là \SI{12.0}{\milli\liter}. Tính nồng độ mol của \ce{NaOH}.
\dapso{\SI{0.0600}{\mole\per\liter}.}
\begin{loigiaiBT}
\[
n_{\ce{HCl}}=0.0120\times0.0500=\SI{6.00e-4}{\mole}.
\]
Theo phản ứng \ce{HCl + NaOH -> NaCl + H2O}, số mol \ce{NaOH} bằng số mol \ce{HCl}:
\[
C_{\ce{NaOH}}=\frac{6.00\times10^{-4}}{0.0100}=\SI{0.0600}{\mole\per\liter}.
\]

\textbf{Vậy,} nồng độ mol của dung dịch \ce{NaOH} là \SI{0.0600}{\mole\per\liter}.
\end{loigiaiBT}
\end{tuluyen}
